% BHU PYQ -- Auto-generated & error-corrected
\documentclass[11pt,a4paper]{article}

\usepackage[a4paper,top=2.5cm,bottom=2.5cm,left=2.8cm,right=2.8cm,headheight=14pt]{geometry}
\usepackage{amsmath,amssymb}
\usepackage[T1]{fontenc}
\usepackage[utf8]{inputenc}
\usepackage{lmodern}
\usepackage{microtype}
\usepackage{fancyhdr}
\usepackage{enumitem}
\usepackage{booktabs}
\usepackage{hyperref}
\usepackage{lastpage}
\usepackage{parskip}

\hypersetup{colorlinks=true,urlcolor=black,linkcolor=black,
  pdftitle={STA/STB-601: Numerical Analysis -- BHU PYQ},pdfauthor={Banaras Hindu University}}

\pagestyle{fancy}\fancyhf{}
\renewcommand{\headrulewidth}{0.4pt}\renewcommand{\footrulewidth}{0.4pt}
\fancyhead[L]{\small\textit{Banaras Hindu University}}
\fancyhead[R]{\small\textit{STA/STB-601: Numerical Analysis}}
\fancyfoot[C]{\small Page \thepage\ of \pageref{LastPage}}

\newlist{parts}{enumerate}{2}
\setlist[parts,1]{label=(\alph*),leftmargin=*,itemsep=5pt,topsep=3pt}
\setlist[parts,2]{label=(\roman*),leftmargin=*,itemsep=3pt,topsep=2pt}
\newcommand{\pts}[1]{\hfill\textit{[#1]}}

\begin{document}

\begin{center}
  {\large\bfseries BANARAS HINDU UNIVERSITY}\\[2pt]
  {\normalsize B.A./B.Sc. (Hons.) Semester VI Examination, 2013-14}\\[6pt]
  \rule{0.8\linewidth}{0.5pt}\\[6pt]
  {\large\bfseries Paper No. STA/STB-601: Numerical Analysis}\\[4pt]
  {\normalsize Time: 3 Hours\hfill Maximum Marks: 70}
\end{center}
\rule{\linewidth}{0.4pt}
\smallskip
\noindent\textit{Instructions: Answer any \textbf{five} questions. All questions carry equal marks. Symbols have their usual meanings.}
\bigskip

\medskip\noindent\textbf{Question 1.}
\begin{parts}
  \item Prove the following identities (symbols have their usual meaning):
  \begin{parts}
    \item $\Delta \cdot E^{-1} = \nabla$
    \item $(\Delta - \nabla)(\Delta + \nabla) = \Delta^2 - \nabla^2$
  \end{parts}\pts{7}
  \item State and prove the fundamental theorem of difference calculus. \pts{7}
\end{parts}

\medskip\rule{\linewidth}{0.2pt}

\medskip\noindent\textbf{Question 2.}
\begin{parts}
  \item Show that $\Delta^n(x^{(n)}) = n!$ and hence find the value of $\Delta^6(x^6)$. \pts{7}
  \item Obtain the missing term $y_2$ in the following table:

  \begin{center}
  \begin{tabular}{c|ccccc}
  \toprule
  $x$ & 0 & 1 & 2 & 3 & 4 \\
  \midrule
  $y$ & 1 & 3 & -- & 13 & 21 \\
  \bottomrule
  \end{tabular}
  \end{center}\pts{7}
\end{parts}

\medskip\rule{\linewidth}{0.2pt}

\medskip\noindent\textbf{Question 3.}
\begin{parts}
  \item State and prove the Newton--Gregory forward interpolation formula. \pts{7}
  \item Using Stirling's interpolation formula, find the value of $y_{35}$ given that $y_{20} = 512$, $y_{30} = 439$, $y_{40} = 346$, $y_{50} = 243$. \pts{7}
\end{parts}

\medskip\rule{\linewidth}{0.2pt}

\medskip\noindent\textbf{Question 4.}
\begin{parts}
  \item State and prove Newton's divided difference interpolation formula for unequal intervals. \pts{7}
  \item Using the Lagrange interpolation formula, find the polynomial form of the function $f(x)$ tabulated below:

  \begin{center}
  \begin{tabular}{c|cccc}
  \toprule
  $x$ & 0 & 1 & 4 & 5 \\
  \midrule
  $f(x)$ & 4 & 3 & 24 & 39 \\
  \bottomrule
  \end{tabular}
  \end{center}\pts{7}
\end{parts}

\medskip\rule{\linewidth}{0.2pt}

\medskip\noindent\textbf{Question 5.}
\begin{parts}
  \item What is inverse interpolation? Tabulate $y = x^{1/3}$ for $x = 2, 3, 4, 5$ and calculate the cube root of 10 correct to three decimal places. \pts{7}
  \item From the following table, compute $\dfrac{dy}{dx}$ and $\dfrac{d^2y}{dx^2}$ at $x = 1$:

  \begin{center}
  \begin{tabular}{c|ccccc}
  \toprule
  $x$ & 1.0 & 1.5 & 2.0 & 2.5 & 3.0 \\
  \midrule
  $y$ & 2.7183 & 4.4817 & 7.3891 & 12.182 & 20.086 \\
  \bottomrule
  \end{tabular}
  \end{center}\pts{7}
\end{parts}

\medskip\rule{\linewidth}{0.2pt}

\medskip\noindent\textbf{Question 6.}
\begin{parts}
  \item State and prove Simpson's three-eighths ($3/8$) rule of numerical integration. \pts{7}
  \item Calculate the value of $\pi$ from $\displaystyle\int_0^1 \frac{dx}{1+x^2}$, taking six equal intervals and using Simpson's one-third rule. \pts{7}
\end{parts}

\medskip\rule{\linewidth}{0.2pt}

\medskip\noindent\textbf{Question 7.}
\begin{parts}
  \item Sum the following series to the first $n$ terms:
  \[
    1, \; 12, \; 36, \; 80, \; 150, \; 252, \; \ldots
  \]\pts{7}
  \item Find the sum of the first $n$ terms of a series whose $k$-th term $T_k$ is given by $T_k = 2^k(3 + k)$. \pts{7}
\end{parts}

\medskip\rule{\linewidth}{0.2pt}

\medskip\noindent\textbf{Question 8.}
\begin{parts}
  \item Using the Runge--Kutta method, approximate $y$ when $x = 0.2$, given that $y = 1$ when $x = 0$ and $\dfrac{dy}{dx} = x + y$. \pts{7}
  \item Solve $\dfrac{dy}{dx} = 1 - y$, $y(0) = 0$, in the range $0 \leq x \leq 0.3$ using the modified Euler's method with step size $h = 0.1$. \pts{7}
\end{parts}

\vfill
\rule{\linewidth}{0.4pt}
\begin{center}\small
  \href{https://bhu-exam.vercel.app/}{Click here to visit our website}\\[4pt]
  \href{https://me-aryan.vercel.app/}{Click here to visit our contributor}\\[4pt]
  \href{https://quashan-me.vercel.app/}{Click here to visit our contributor}
\end{center}

\end{document}